% ===========================================================================
\chapter{A conversation as data}\label{ch:conversation}

Before the agent can talk to a model, we need a way to write the
conversation down. This chapter builds the data: messages, the tool calls
a model makes, the results we send back, and the model's replies. It uses
the shapes of Claude's Messages API, so that in Chapter~\ref{ch:claude}
sending a conversation is only a matter of turning it into JSON.

\section{What the model sees}

A language model has no memory between requests. Each request carries the
whole conversation so far, and the model writes the next message. A
conversation is a list of messages, each from one of two \emph{roles}:
\code{"user"} (us) and \code{"assistant"} (the model). They take turns,
and the first message is the user's. Each message holds a list of
\emph{content blocks}. The simplest is text. Here is the start of a
conversation, as JSON:

\begin{json}
[{"role": "user",
  "content": [{"type": "text", "text": "Some of the checks fail. Find out why and fix it."}]}]
\end{json}

The model's reply is another message. Along with the request, we tell the
model which tools it may use (Chapter~\ref{ch:tools}), and if it wants one,
its reply contains a \code{tool_use} block. The block names the tool,
gives its input as a JSON object, and has an \code{id}:

\begin{json}
{"role": "assistant",
 "content": [{"type": "text", "text": "I'll start by running the checks."},
             {"type": "tool_use", "id": "toolu_01", "name": "run_command",
              "input": {"command": "nex checks.nex"}}],
 "stop_reason": "tool_use"}
\end{json}

The \code{stop_reason} says why the model stopped writing. Here it is
\code{"tool_use"}: it is waiting for the result. The result goes in the
next user message, as a \code{tool_result} block quoting the call's
\code{id}:

\begin{json}
{"role": "user",
 "content": [{"type": "tool_result", "tool_use_id": "toolu_01",
              "content": "exit code 1\nPASS a loan is due 14 days ...", "is_error": false}]}
\end{json}

A reply may ask for several tools at once, each with its own \code{id}.
Every one of them must be answered, all in the next message. When the
model has finished, its \code{stop_reason} is \code{"end_turn"}.

That is all the structure there is. The rest of this chapter writes it in
Nex.

\section{JSON in Nex}

Nex's \code{data/Json} library reads JSON text into a \code{Json_Value},
which we can take apart with \code{get} (a field of an object), \code{at}
(an element of an array), \code{items} (all the elements), and
\code{as_string}, \code{as_integer} and so on, which check the value's type
as they go. In the other direction, \code{stringify} turns Nex maps, arrays,
strings and numbers into JSON text. Two short functions make both
directions easier to write:

\begin{newfile}{messages.nex}
-- messages.nex: a conversation with the model, as data.
--
-- The Claude API is stateless: every request carries the whole
-- conversation so far. A conversation is a list of messages, alternating
-- between the user and the assistant (the model). Each message holds a
-- list of content blocks: text, a tool call the model makes ("tool_use"),
-- or the result of one ("tool_result").

intern data/Json

-- A JSON text read back as a Json_Value: json_value('{"path": "a.nex"}').
function json_value(text: String): Json_Value do
  result := create Json.make().read(text)
end

-- Any value built from Nex maps and arrays, as JSON text.
function json_text(value: Any): String do
  result := create Json.make().stringify(value)
end
\end{newfile}

The example in the comment uses a \emph{raw} string, in single quotes.
Nothing inside one is special: no backslash escapes, and double quotes are
just characters. That makes raw strings the easy way to write JSON in Nex
code, and we will use them a lot.

\section{Tool calls and their results}

A tool call is three things: the \code{id}, the tool's name, and its
input. The input stays a \code{Json_Value}, because each tool expects
different arguments, and each tool will take apart its own.

\begin{append}{messages.nex}
-- The model asking for one tool to be run.
class Tool_Call
create
  make(id, name: String, input: Json_Value) do
    this.id := id
    this.name := name
    this.input := input
  end

feature
  id: String          -- chosen by the model; the result must quote it
  name: String        -- which tool
  input: Json_Value   -- the tool's arguments, a JSON object
end
\end{append}

The constructor's parameters have the same names as the fields, so it
writes \code{this.id} for the field. (Writing \code{id := id} would assign
the parameter to itself, and leave the field empty.)

A result is the text a tool gave back, and whether it failed. That second
part matters more than it looks. A failed tool is not the end of the
world: the model reads the error, just as a programmer would, and tries
something else. But it needs to know that it \emph{was} an error. The
\code{to_json} routine builds the block the API wants, as a Nex map:

\begin{append}{messages.nex}
-- What running one tool gave back, ready to send to the model.
class Tool_Result
create
  make(call_id, text: String, is_error: Boolean) do
    this.call_id := call_id
    this.text := text
    this.is_error := is_error
  end

feature
  call_id: String
  text: String
  is_error: Boolean   -- tells the model the tool failed, and why

  to_json(): Map[String, Any] do
    result := {"type": "tool_result", "tool_use_id": call_id,
      "content": text, "is_error": is_error}
  end
end
\end{append}

\code{Map[String, Any]} is a map from strings to values of any type: here,
strings and a boolean. That is what a JSON object is.

\section{Replies}

A reply arrives from the API as a JSON object. The obvious thing to do
would be to take it apart into Nex objects (some text, some tool calls)
and build the JSON again when we send the conversation back. We will not
do that, and the reason is worth knowing. Models like Claude can
\emph{think} before they answer, and a reply then contains
\code{thinking} blocks too, with a signature that the API checks when they
come back. Any change to them, or to the blocks around them, and they are
no longer valid. The safe rule, and the one we follow, is: \emph{a reply
goes back to the API exactly as it came}. So a \code{Reply} keeps the whole
response, and reads things out of it when asked:

\begin{append}{messages.nex}
-- One reply from the model, kept exactly as the API sent it.
--
-- Its content blocks go back to the API unchanged in every later
-- request. Some of them (the model's "thinking") are only valid if
-- nothing about them changes, so a reply is never rebuilt from the
-- pieces this class reads out of it.
class Reply
create
  -- A reply from the body of a Messages API response.
  from_json(response: Json_Value) do
    this.response := response
  end

feature
  response: Json_Value

  -- The content blocks, a JSON array.
  content(): Json_Value do
    result := response.get("content")
  end

  -- Why the model stopped: "end_turn" (it has finished), "tool_use" (it
  -- wants tools run), "max_tokens" (it ran out of room), "refusal", ...
  stop_reason(): String do
    result := response.get("stop_reason").as_string
  end
\end{append}

The text of a reply is its text blocks, joined. The tool calls are its
\code{tool_use} blocks, in order. Blocks of any other type, thinking
included, are simply skipped: they are for the API, not for us.

\begin{continued}{messages.nex}
  -- The text of the reply: its text blocks, one per paragraph.
  text(): String do
    let parts: Array[String] := []
    across content().items as block do
      if block.get("type").as_string = "text" then
        parts.add(block.get("text").as_string)
      end
    end
    result := "\n\n".join(parts)
  end

  -- The tools the model wants run, in the order it asked for them.
  tool_calls(): Array[Tool_Call] do
    across content().items as block do
      if block.get("type").as_string = "tool_use" then
        result.add(create Tool_Call.make(block.get("id").as_string,
          block.get("name").as_string, block.get("input")))
      end
    end
  end
\end{continued}

\texttt{"\textbackslash n\textbackslash n".join(parts)} joins the parts with
a blank line between them.

A response from the API also says how many tokens the request used, in its
\code{usage} field. We will want that in Chapter~\ref{ch:long}, to see how
long the conversation is getting. The API can keep a request's start in a
\emph{prompt cache}, so that the next request, which starts the same way,
is cheaper; then the tokens read are split into three counts (fresh, read
from the cache, and written to it), and the length of the conversation is
their sum:

\begin{continued}{messages.nex}
  -- How many tokens the model read: the whole request, cached or not.
  -- This is how full the context window is.
  context_tokens(): Integer do
    result := usage("input_tokens") + usage("cache_read_input_tokens")
      + usage("cache_creation_input_tokens")
  end

  -- How many of those were read from the prompt cache, at a tenth of
  -- the price.
  cached_tokens(): Integer do
    result := usage("cache_read_input_tokens")
  end

  -- How many tokens the model wrote.
  output_tokens(): Integer do
    result := usage("output_tokens")
  end

  private feature
    usage(field: String): Integer do
      if response.has("usage") and response.get("usage").has(field)
        and response.get("usage").get(field).is_number then
        result := response.get("usage").get(field).as_integer
      end
    end
end
\end{continued}

A \code{private feature} section is for the class's own use: other code
can't call \code{usage}. If a count is missing, as it will be in the
replies we write by hand, it is 0.

\section{Replies written by hand}

Until Chapter~\ref{ch:claude} the model will be played by a script, so we
need replies we can write ourselves. \code{scripted_reply} builds the same
JSON the API would send: an optional text block, and an optional tool call,
whose input is given as JSON text. It then turns it into text and reads it
back, so that a reply we write and a reply from the API are made of exactly
the same kind of values.

\begin{append}{messages.nex}
-- A reply written by hand, for a pretend model: `text` and, if `tool`
-- is not "", one call of `tool` with `input` (a JSON object as text).
function scripted_reply(text, tool, input: String): Reply do
  let blocks: Array[Any] := []
  if text /= "" then
    blocks.add({"type": "text", "text": text})
  end
  let stop: String := "end_turn"
  if tool /= "" then
    blocks.add({"type": "tool_use", "id": "call_" + tool + "_" + input.hash.abs.to_string,
      "name": tool, "input": json_value(input).raw()})
    stop := "tool_use"
  end
  let response: Map[String, Any] := {"type": "message", "role": "assistant",
    "content": blocks, "stop_reason": stop}
  result := create Reply.from_json(json_value(json_text(response)))
end

-- A reply that only talks.
function say(text: String): Reply do
  result := scripted_reply(text, "", "")
end

-- A reply that says `text` and asks for `tool` to be run on `input`.
function use_tool(text, tool, input: String): Reply do
  result := scripted_reply(text, tool, input)
end
\end{append}

\code{raw()} gives back the plain Nex value inside a \code{Json_Value}, to
put inside another one. The call's \code{id} is made from the tool and a
hash of its input, so it is the same every time the script runs.

\section{The conversation}

Now the conversation itself. It keeps, for each message, its role and its
content blocks. It also keeps the ids of the tool calls in the last reply
that haven't been answered yet, so that it can refuse to go on until they
are.

\begin{append}{messages.nex}
-- The conversation so far.
class Conversation
create
  make() do
  end

feature
  roles: Array[String]           -- "user" or "assistant", one per message
  contents: Array[Array[Any]]    -- each message's content blocks
  waiting: Array[String]         -- ids of tool calls not yet answered

  length(): Integer do
    result := roles.length
  end

  -- Whether the conversation ends with the user, so it is the model's
  -- turn to speak.
  models_turn(): Boolean do
    result := not roles.is_empty and roles.get(roles.length - 1) = "user"
      and waiting.is_empty
  end
\end{append}

There are three ways the conversation grows: the user says something, the
model replies, or we send back tool results. Each comes with a contract
that keeps the conversation the way the API wants it.

When the user speaks right after a user message (as happens when we send
tool results and then have something to add), the new text joins that
message instead of starting another one. That way the roles always take
turns. But the user may not speak at all while a tool call is waiting for
its answer: the API would reject the conversation.

\begin{continued}{messages.nex}
  -- Add something the user says. If the conversation already ends with
  -- the user (with tool results, say), it joins that message, so that
  -- user and assistant still take turns.
  add_user_text(text: String)
  require
    something_said: text /= ""
    no_calls_waiting: waiting.is_empty
  do
    let block: Map[String, Any] := {"type": "text", "text": text}
    add_user_block(block)
  ensure
    models_turn: models_turn()
  end
\end{continued}

A reply goes in exactly as it came, block by block, and its tool calls
start waiting:

\begin{continued}{messages.nex}
  -- Add a reply from the model, exactly as it came.
  add_reply(reply: Reply)
  require
    models_turn: models_turn()
  do
    roles.add("assistant")
    let blocks: Array[Any] := []
    across reply.content().items as block do
      blocks.add(block.raw())
    end
    contents.add(blocks)
    across reply.tool_calls() as call do
      waiting.add(call.id)
    end
  end
\end{continued}

Tool results must answer the waiting calls exactly: one result for each,
no more, no fewer, all in one message.

\begin{continued}{messages.nex}
  -- Answer the tool calls of the last reply: one result for every call,
  -- all in one message.
  add_tool_results(results: Array[Tool_Result])
  require
    answers_every_call: answers_exactly(results)
  do
    across results as r do
      add_user_block(r.to_json())
    end
    waiting := []
  ensure
    models_turn: models_turn()
  end

  -- Whether `results` answer the calls that are waiting, each once.
  answers_exactly(results: Array[Tool_Result]): Boolean do
    result := results.length = waiting.length and not waiting.is_empty
    across results as r do
      if not waiting.contains(r.call_id) then
        result := false
      end
    end
  end
\end{continued}

These preconditions turn mistakes that the API would report much later,
with a confusing error, into an immediate stop at the line that made them.

The remaining routines read the conversation: the last user message (we
will need it to record sessions, in Chapter~\ref{ch:replay}), the whole
conversation as the API wants it, and whether the roles take turns.

\begin{continued}{messages.nex}
  -- The last message the user sent: what a request adds to the ones
  -- before it.
  last_user_message(): Array[Any]
  require
    there_is_one: roles.contains("user")
  do
    let i: Integer := roles.length - 1
    from until roles.get(i) = "user" do
      i := i - 1
    end
    result := contents.get(i)
  end

  -- The messages as the API wants them.
  to_json(): Array[Any] do
    let i: Integer := 0
    across roles as role do
      result.add({"role": role, "content": contents.get(i)})
      i := i + 1
    end
  end

  -- Whether no two messages in a row come from the same side.
  alternates(): Boolean do
    result := true
    let i: Integer := 1
    from until i >= roles.length do
      if roles.get(i) = roles.get(i - 1) then
        result := false
      end
      i := i + 1
    end
  end
\end{continued}

Last, the private routine that adds a block to the user's side, and the
class \emph{invariant}: what must be true of every conversation, before
and after every call of its routines. Nex checks it each time.

\begin{continued}{messages.nex}
  private feature
    add_user_block(block: Map[String, Any]) do
      if roles.is_empty or roles.get(roles.length - 1) = "assistant" then
        roles.add("user")
        contents.add([])
      end
      contents.get(contents.length - 1).add(block)
    end

invariant
  one_content_per_role: roles.length = contents.length
  starts_with_user: roles.is_empty or roles.get(0) = "user"
  takes_turns: alternates()
end
\end{continued}

Let's build a conversation, and see it as the API will:

\begin{tryit}{try_conversation.nex}
intern Messages

let c: Conversation := create Conversation.make()
c.add_user_text("Some of the checks fail. Find out why and fix it.")
let reply: Reply := use_tool("I'll start by running the checks.", "run_command",
  '{"command": "nex checks.nex"}')
c.add_reply(reply)
console.print_line("waiting for: " + c.waiting.to_string)

let call: Tool_Call := reply.tool_calls().get(0)
console.print_line(call.name + " " + call.input.get("command").as_string)
c.add_tool_results([create Tool_Result.make(call.id, "exit code 1\n...", false)])
console.print_line("the model's turn: " + c.models_turn().to_string)
across c.to_json() as message do
  console.print_line(json_text(message))
end

let refused: String := ""
do
  c.add_reply(say("Hello again."))
  c.add_reply(say("And again."))
rescue
  refused := "" + exception
end
console.print_line(refused)
\end{tryit}
\tryoutput{try_conversation}

The last part breaks a rule on purpose: two replies in a row. The second
\code{add_reply} stops with a precondition violation, which the
\code{rescue} clause catches. (In a \code{rescue}, \code{exception} is
what went wrong; adding it to a string gives its message.)
